% 原书第62--64页；2.4全部四分目.

\section{Eichler--Shimura关系}\label{sec:modular-2-4}

为完成与 $f\in\mathcal S_2(\Gamma_0(N))$ 相联系的表示 $\rho_f$ 的构造，我们必须比较 $T(p)$ 与Frobenius代换 $\sigma_p$ 的作用，至少对除有限多个以外的所有素数 $p$ 都应如此.目标是下面的定理.



\subsection{Eichler--Shimura关系定理}\label{sec:modular-2-4-1}

\begin{theorem*}

若素数 $p$ 不整除 $\ell N$，则 $\operatorname{Gal}(\overline{\mathbb Q}/\mathbb Q)$ 在群 $\operatorname{Pic}^0(X_0(N)_{\overline{\mathbb Q}})[\ell^m]$ 上的作用在 $p$ 处非分歧，且对 $p$ 处任意Frobenius代换 $\sigma_p$，在 $\operatorname{End}(\operatorname{Pic}^0(X_0(N)_{\overline{\mathbb Q}})[\ell^m])$ 中有等式
\begin{align*}
T(p)=\sigma_p+p\sigma_p^{-1}.
\end{align*}
\end{theorem*}

证明使用模曲线 $X_0(N)$ 的模 $p$ 约化，以及关于Jacobi簇的一些一般事实.



\subsection{\texorpdfstring{$\mathbb F_p$}{Fₚ} 上的曲线}\label{sec:modular-2-4-2}

设 $C$ 为 $\mathbb F_p$ 上的紧合光滑曲线，并且其标量扩张 $C_{\overline{\mathbb F}_p}$ 仍为紧合光滑曲线（唯一可能不满足的是连通性；换言之，我们要求 $\mathbb F_p$ 在 $\mathbb F_p(C)$ 中的代数闭包就是 $\mathbb F_p$）.如前所述，可以考虑 $\operatorname{Pic}^0(C_{\overline{\mathbb F}_p})$，它带有 $\operatorname{Gal}(\overline{\mathbb F}_p/\mathbb F_p)$ 的作用，并与 $C$ 的Jacobi簇的 $\overline{\mathbb F}_p$ 值点相认同.以 $\sigma_p$ 表示 $\operatorname{Gal}(\overline{\mathbb F}_p/\mathbb F_p)$ 的通常生成元 $x\mapsto x^p$.



域 $\mathbb F_p(C)$ 上有Frobenius自同态 $F^*:f\mapsto f^p$.扩张标量后，它诱导环 $\overline{\mathbb F}_p\otimes_{\mathbb F_p}\mathbb F_p(C)$ 的自同态 $\mathrm{id}\otimes F^*$；后者延拓至其分式域 $\overline{\mathbb F}_p(C)$，称为``几何Frobenius''.注意，后者不再由 $f\mapsto f^p$（``绝对Frobenius''）给出，因为它是 $\overline{\mathbb F}_p$ 线性的.实际上，我们有分解
\begin{align*}
(\mathrm{id}\otimes F^*)(\sigma_p\otimes\mathrm{id})=\overline{\mathbb F}_p(C)\text{ 上的绝对Frobenius}.
\end{align*}
用 $F:C_{\overline{\mathbb F}_p}\longrightarrow C_{\overline{\mathbb F}_p}$ 表示由域与曲线的对偶性得到的``几何Frobenius自同态''.在 $C_{\overline{\mathbb F}_p}$ 的闭点上，绝对Frobenius的作用平凡，故相对Frobenius的作用与Galois群的生成元 $\sigma_p$ 相同.具体而言，若将 $C$ 嵌入 $\mathbb P^r_{\mathbb F_p}$，则 $F([x_0,\ldots,x_r])=[x_0^p,\ldots,x_r^p]$.然而，尽管 $F$ 诱导闭点的一个置换，它并非自同构：它在每点都有分歧（甚至是不可分的），次数为 $p$.



与任何非常值代数曲线态射一样，自同态 $F$ 诱导两个自同态
\begin{align*}
F_*\quad\text{和}\quad F^*\in\operatorname{End}\bigl(\operatorname{Pic}^0(C_{\overline{\mathbb F}_p})\bigr).
\end{align*}
由前面的讨论，可以将其作用与Galois作用比较：
\begin{align*}
F_*=\sigma_p,\qquad F^*=p\sigma_p^{-1}.
\end{align*}
\subsection{曲线的约化}\label{sec:modular-2-4-3}

设 $C$ 为 $\mathbb Q$ 上的光滑射影曲线，$p$ 为素数.$C$ 在局部化 $\mathbb Z_{(p)}$ 上的模型 $\mathcal C$，是 $\mathbb Z_{(p)}$ 上满足
\begin{align*}
\mathcal C_{\mathbb Q}:=\mathcal C\times_{\operatorname{Spec}(\mathbb Z_{(p)})}\operatorname{Spec}(\mathbb Q)\simeq C
\end{align*}
的平坦概形.若其约化
\begin{align*}
\mathcal C_{\mathbb F_p}:=\mathcal C\times_{\operatorname{Spec}(\mathbb Z_{(p)})}\operatorname{Spec}(\mathbb F_p)
\end{align*}
为光滑曲线，就称它为好模型.具体而言，为了找到模型，可在 $\mathbb P^r$ 中用系数属于 $\mathbb Z_{(p)}$ 的齐次多项式定义 $C$；约化 $\mathcal C_{\mathbb F_p}$ 便由这些多项式模 $p$ 后定义.但寻找好模型更加困难，有时好模型并不存在.若存在好模型，就说 $C$ 有\newterm{modular-good-reduction}{好约化}{good reduction}.



在这种情形，约化映射 $\mathbb P^r(\mathbb Q)\longrightarrow\mathbb P^r(\mathbb F_p)$（注意，齐次坐标中可消去分母）诱导从 $C$ 的 $\mathbb Q$ 值点到 $\mathcal C_{\mathbb F_p}$ 的 $\mathbb F_p$ 值点的映射.更一般地，若 $\mathfrak p$ 为 $\overline{\mathbb Z}$ 中位于 $p$ 上方的极大理想，则模 $\mathfrak p$ 约化诱导映射
\begin{equation*}
\begin{tikzcd}
C(\overline{\mathbb Q}) \arrow[r] & \mathcal C_{\mathbb F_p}(\overline{\mathbb F}_p).
\end{tikzcd}
\end{equation*}
此映射线性延拓到除子上.可以证明以下性质（仍在好约化的假设下）：

\begin{enumerate}[label=\roman*)]

\item 模 $\mathfrak p$ 约化映射诱导满同态
\begin{equation*}
\begin{tikzcd}
\operatorname{Pic}^0(C_{\overline{\mathbb Q}}) \arrow[r] & \operatorname{Pic}^0(C_{\overline{\mathbb F}_p}).
\end{tikzcd}
\end{equation*}
\item 对 $\ell\ne p$，它诱导同构
\begin{equation*}
\begin{tikzcd}
\operatorname{Pic}^0(C_{\overline{\mathbb Q}})[\ell^m] \arrow[r,"\sim"] & \operatorname{Pic}^0(C_{\overline{\mathbb F}_p})[\ell^m].
\end{tikzcd}
\end{equation*}
\end{enumerate}



\subsection{模曲线 \texorpdfstring{$X_0(N)_{\mathbb Q}$}{X₀(N) 在Q上} 的约化}\label{sec:modular-2-4-4}

可以证明，$X_0(N)$ 在每个 $p\nmid N$ 处都有好约化.从已经得到的``模''方程 $f_N(X,Y)$ 看出这一点似乎很困难.更适合的方法，是使用 $X_0(N)$ 参数化带有 $p$ 阶子群的椭圆曲线的模解释.\footnote{译注：原书写作``$p$ 阶子群''；通常 $X_0(N)$ 参数化带有 $N$ 阶循环子群的椭圆曲线.此处按原文保留.}由于没有时间讨论这一点，我们只指出，$\mathbb Q$ 上的曲线在除有限多个素数之外总有好约化.设 $p$ 是这样的素数.则约化 $X_0(N)_{\mathbb F_p}$ 的函数域，可以在 $\mathbb F_p(j)$ 上由同一方程 $f_N(X,Y)$ 模 $p$ 后定义，即
\begin{align*}
\mathbb F_p(X_0(N)_{\mathbb F_p})=\mathbb F_p(j)[Y]/(f_N(j,Y)).
\end{align*}
我们希望用 $F_*$ 和 $F^*$ 来理解 $T(p)$.对应 $T(p)$ 涉及 $X_0(N,p)$，而后者遗憾地在 $p$ 处没有好约化.不过，联系 $j(z)$ 与 $j(p^{-1}z)$ 的方程 $f_p$ 具有一个有趣的约化，由定理~\ref{sec:modular-2-3-3} 给出.



为简化讨论，令 $N=1$，尝试理解对应
\begin{equation*}
\begin{tikzcd}
X(1)_{\mathbb F_p} & X_0(1,p)_{\mathbb F_p} \arrow[l,"\pi^p"'] \arrow[r,"\pi"] & X(1)_{\mathbb F_p},
\end{tikzcd}
\end{equation*}
至少在普遍点附近，即在全分式环上理解它.根据定理~\ref{sec:modular-2-3-3}，``对偶地''有
\begin{equation*}
\begin{tikzcd}[column sep=small]
\mathbb F_p(j) \arrow[r] & {\mathbb F_p(j)[Y]/\bigl((Y^p-j)(Y-j^p)\bigr)} & \mathbb F_p(j) \arrow[l].
\end{tikzcd}
\end{equation*}
也就是\footnote{译注：原书右侧因子写为 $\mathbb F_p(j)/(Y-j^p)$，缺少多项式变量 $[Y]$；这里保留原式.}
\begin{equation*}
\begin{tikzcd}[column sep=small]
\mathbb F_p(j) \arrow[r] & {\mathbb F_p(j)[Y]/(Y^p-j)\times\mathbb F_p(j)/(Y-j^p)} & \mathbb F_p(j) \arrow[l],
\end{tikzcd}
\end{equation*}
最后为
\begin{equation*}
\begin{tikzcd}[column sep=small]
\mathbb F_p(j) \arrow[r] & {\mathbb F_p(j^{1/p})\times\mathbb F_p(j)} & \mathbb F_p(j) \arrow[l].
\end{tikzcd}
\end{equation*}
在此图中，右边箭头将 $j$ 映为 $(j,j)$，左边箭头则将 $j$ 映为 $(Y,Y)=(j^{1/p},j^p)$.对于 $X_0(N)$，类似的图表明，在普遍点附近，对应具有如下形式：
\begin{equation*}
\begin{tikzcd}[column sep=small]
X_0(N)_{\mathbb F_p} & {X_0(N)_{\mathbb F_p}\bigcup X_0(N)_{\mathbb F_p}} \arrow[l,"\mathrm{id}\cup F"'] \arrow[r,"F\cup\mathrm{id}"] & X_0(N)_{\mathbb F_p}.
\end{tikzcd}
\end{equation*}
这里``在普遍点附近''是说，在 $X_0(N)_{\mathbb F_p}$ 的某个开集之上，情形确实如上所述（为不交并）.但``整个'' $X_0(N,p)$ 的约化，是两个 $X_0(N)_{\mathbb F_p}$ 副本的非不交并.



我们承认，这一普遍点附近的情形足以说明，对应 $T(p)$ 与 $(F\cup\mathrm{id})_*\circ(\mathrm{id}\cup F)^*=F+F^*$ 相同；\footnote{译注：原书末式写作 $F+F^*$；结合本节前文，作为 $\operatorname{Pic}^0$ 的自同态之和应读作 $F_*+F^*$.}这就完成了Eichler--Shimura关系的证明（更确切地说，是解释）.

